\documentclass[10pt,reqno]{amsart}

\usepackage[T1]{fontenc}
\usepackage[utf8]{inputenc}
\usepackage{newtxtext,newtxmath}
\usepackage{microtype}
\usepackage{mathtools}
\usepackage{booktabs}
\usepackage{enumitem}
\usepackage{xcolor}
\usepackage{hyperref}
\usepackage[a4paper,margin=28mm,headheight=15pt]{geometry}

\definecolor{LinkBlue}{HTML}{174A7E}
\hypersetup{
  colorlinks=true,
  linkcolor=LinkBlue,
  citecolor=LinkBlue,
  urlcolor=LinkBlue
}

\setlength{\parskip}{2pt}
\setlist[itemize]{leftmargin=6mm,itemsep=1.5pt,topsep=3pt}
\setlist[enumerate]{leftmargin=7mm,itemsep=2pt,topsep=3pt}
\renewcommand{\arraystretch}{1.15}
\emergencystretch=2em

\newtheorem{theorem}{Theorem}[section]
\newtheorem{proposition}[theorem]{Proposition}
\newtheorem{lemma}[theorem]{Lemma}
\newtheorem{corollary}[theorem]{Corollary}
\theoremstyle{definition}
\newtheorem{definition}[theorem]{Definition}
\newtheorem{remark}[theorem]{Remark}

\title[Regularity-Premise Obstruction in Fu's Theorem 2.8]
{A Local-Finiteness Obstruction to the Regularity Premise in Fu's Theorem 2.8}

\author{Argus Open Research Project}
\address{Argus Open Research Project}
\date{August 29, 2026}

\subjclass[2020]{Primary 03C64; Secondary 14P15, 26E10}
\keywords{o-minimal structures, Abel equation, trans-exponential growth, locally finite triangulation}

\begin{document}

\begin{abstract}
In a 2026 manuscript investigating trans-exponential o-minimal expansions of the real ordered field, Theorem~2.8 assumes that for every allowed polynomial-exponential system $P$, the regular parameter locus $R\subset G\times[-1,1]^{2n+2}$ admits a triangulation by simplices $\{\Delta_\alpha\}$ that is locally finite in the full ambient parameter space $G\times[-1,1]^{2n+2}$. We prove that for the allowed system $n=1,\ k=u=0,\ P(x,l_1,l_2,\epsilon,\delta)=x^2$, the identity-coordinate slice of $R$ in $[-1,1]$ is $[-1,1]\setminus\{0\}$, which is not closed. Because the union of any locally finite family of closed subsets in a topological space is closed, no locally finite triangulation of $R$ by closed simplices exists in $G\times[-1,1]^{2n+2}$. If the simplices $\Delta_\alpha$ are instead taken to be open simplex interiors, then the closedness step in the subsequent Claim~2.9 fails. The topological and logical core is verified in Lean~4.
\end{abstract}

\maketitle

\section{Main Obstruction Theorem}

Let $n=1$ and $k=u=0$, and consider the allowed polynomial system
\[
  P(x,l_1,l_2,\epsilon,\delta) \;=\; x^2.
\]
For $(A,\eta,l,\epsilon,\delta)\in G\times[-1,1]^4$, the regular parameter set is defined by
\[
  R \;=\; \bigl\{(A,\eta,l,\epsilon,\delta)\in G\times[-1,1]^4 : P\circ A - \eta\text{ has nonsingular zeros only}\bigr\}.
\]
Fixing the identity affine transformation $A=\mathrm{id}$ and auxiliary parameters $(l_1,l_2,\epsilon,\delta)=(0,0,0,0)$, the univariate polynomial $x^2-\eta$ has a singular zero at $x=0$ when $\eta=0$ and only nonsingular real zeros (or no real zeros) when $\eta\in[-1,1]\setminus\{0\}$.

\begin{theorem}\label{thm:main}
For $P(x,l_1,l_2,\epsilon,\delta)=x^2$, the slice of $R$ along the identity coordinate $A=\mathrm{id}$ and $(l_1,l_2,\epsilon,\delta)=0$ is
\[
  R_{\mathrm{id}} \;=\; [-1,1]\setminus\{0\}.
\]
Consequently:
\begin{enumerate}[label=\textup{(\arabic*)}]
  \item $R$ is not closed in $G\times[-1,1]^4$.
  \item There exists no family $\{\Delta_\alpha\}$ of closed simplices with $\bigcup_\alpha\Delta_\alpha = R$ that is locally finite at every point of $G\times[-1,1]^4$.
  \item If $\{\Delta_\alpha\}$ is instead a locally finite family of open simplices with $\bigcup_\alpha\Delta_\alpha = R$, then individual members $\Delta_\alpha$ are not closed in $G\times[-1,1]^4$, obstructing the closed-membership inference in Claim~2.9.
\end{enumerate}
\end{theorem}

\begin{proof}
Since $0\in\overline{[-1,1]\setminus\{0\}}\setminus R_{\mathrm{id}}$, $R_{\mathrm{id}}$ is not closed in $[-1,1]$, and by continuity of the slice inclusion $\eta\mapsto(\mathrm{id},\eta,0,0,0,0)$, $R$ is not closed in $G\times[-1,1]^4$. In any topological space $X$, if $\{F_\alpha\}$ is a locally finite collection of closed sets, then $\bigcup_\alpha F_\alpha$ is closed in $X$; applying this to $X=G\times[-1,1]^4$ proves~(2), and~(3) is immediate.
\end{proof}

\begin{remark}[Scope Boundary]
Theorem~\ref{thm:main} identifies a gap in the stated regularity premise of Theorem~2.8 and its application in Claim~2.9. It does not prove or refute the existence of a trans-exponential o-minimal expansion of $(\mathbb{R},<,+,\cdot,\exp)$.
\end{remark}


\section{Methodological Framework and Verification Invariants}

The mathematical analysis for Result~1465 in the domain of model theory, o-minimal structures, and real analytic geometry follows a strict separation between the exact statement decided and nearby formulations with stronger hypotheses. Specifically:
\begin{enumerate}[label=\textup{(\alph*)}]
  \item \emph{Statement alignment:} Every quantifier, ambient category, normalization convention, and regularity or finiteness hypothesis is fixed against the primary problem source prior to evaluating candidate constructions or obstructions.
  \item \emph{Exactness of algebraic and combinatorial steps:} All finite-dimensional linear-algebraic, polynomial, lattice, or graph-theoretic invariants appearing in the proof above are established in exact rational, integral, or symbolic arithmetic without floating-point approximation.
  \item \emph{Classification:} Within the program taxonomy, this result is categorized as a \emph{proof premise correction} deciding the scope: closed-simplex local-finiteness premise in Theorem 2.8 of arXiv:2604.03477v2 for P(x)=x\textasciicircum\{\}2.
\end{enumerate}

\section{Scope and Open Boundaries}

To prevent over-claiming beyond the proved theorem, we record the explicit boundaries that are \emph{not} decided by the argument above:
\begin{itemize}
  \item existence of a trans-exponential o-minimal expansion of the real ordered field.
\end{itemize}

\begin{thebibliography}{9}
\bibitem{PrimarySource}
Primary problem list and source literature for Result~1465 (model theory, o-minimal structures, and real analytic geometry), archived with exact theorem locators and verification certificates in the accompanying repository package.
\end{thebibliography}

\end{document}
